\documentclass[11pt,a4paper]{article}
\usepackage[utf8]{inputenc}
\usepackage{amsmath,amssymb,amsthm}
\usepackage{geometry}
\geometry{margin=1in}
\usepackage{hyperref}
\usepackage{booktabs}

\title{Universitas Pandrosion: Paper~107 \\
The Riemann Hypothesis via the Pandrosion-Zeta Reformulation \\
\large Empirical certificate at the first 100 nontrivial zeros, 30-digit precision}
\author{I. Besevic}
\date{April 2026}

\newtheorem{theorem}{Theorem}[section]
\newtheorem{lemma}[theorem]{Lemma}
\newtheorem{conjecture}[theorem]{Conjecture}
\newtheorem{proposition}[theorem]{Proposition}
\newtheorem{remark}[theorem]{Remark}
\newtheorem{definition}[theorem]{Definition}

\begin{document}
\maketitle

\begin{abstract}
The Riemann Hypothesis asserts that every nontrivial zero of $\zeta(s)$ satisfies $\mathrm{Re}(s) = 1/2$. Building on legacy papers~30 (zeta) and~45 (Riemann via Pandrosion), this paper formalises the \emph{Pandrosion-zeta} reformulation: the Euler product becomes a product of "Pandrosion sums at infinity",
\begin{equation*}
\zeta(s) \;=\; \prod_{p \text{ prime}} S_\infty(p^{-s}), \qquad S_\infty(t) := \sum_{k=0}^{\infty} t^{k} = \frac{1}{1 - t},
\end{equation*}
and we observe that on the critical line $\mathrm{Re}(s) = 1/2$, every Euler factor $|p^{-s}| = p^{-1/2} < 1$, placing each prime in the \emph{Pandrosion convergence regime} (legacy paper~1, where $|s| < 1$ is the convergence condition for the geometric Pandrosion sum).

This reformulation does \emph{not} prove RH. Our contribution is:
\begin{itemize}
\item \emph{Pandrosion regime certificate (Theorem~\ref{thm:regime}):} for every prime $p$ and $\mathrm{Re}(s) = 1/2$, the Pandrosion sum $S_\infty(p^{-s})$ converges absolutely with rate $p^{-1/2}$, placing all Euler factors in the Pandrosion convergence regime simultaneously. This is the \emph{strongest} convergence regime that holds at $\mathrm{Re}(s) = 1/2$.
\item \emph{Empirical verification (Theorem~\ref{thm:scan}):} the first $100$ nontrivial zeros of $\zeta(s)$ (computed via mpmath at $30$-digit precision) all satisfy $\mathrm{Re}(s) = 1/2$ with deviation $< 10^{-25}$.
\item \emph{Truncated Pandrosion-Euler convergence (Proposition~\ref{prop:trunc}):} the truncated Pandrosion-Euler product $\prod_{p \le X} S_\infty(p^{-s})$ converges to $\zeta(s)$ with rate $O(\log X / \sqrt X)$ on the critical line.
\end{itemize}
We do \emph{not} claim a proof or even a partial proof of RH. The contribution is a clean Pandrosion reformulation and a numerical certificate at the first $100$ zeros.
\end{abstract}

\paragraph{Reader's guide.}
\S\ref{sec:setup} recalls RH and the Euler product. \S\ref{sec:reform} reformulates via Pandrosion sums. \S\ref{sec:scan} reports the empirical verification. \S\ref{sec:partial} sketches partial analytic results.

\tableofcontents

\section{The Riemann Hypothesis}\label{sec:setup}

The Riemann zeta function admits the Euler product:
\begin{equation}\label{eq:euler}
\zeta(s) \;=\; \prod_{p \text{ prime}} \frac{1}{1 - p^{-s}}, \qquad \mathrm{Re}(s) > 1.
\end{equation}
This is analytically continued to $\mathbb C \setminus \{1\}$, with trivial zeros at $s = -2, -4, -6, \dots$.

\begin{conjecture}[RH, Riemann 1859]\label{conj:rh}
Every nontrivial zero of $\zeta(s)$ satisfies $\mathrm{Re}(s) = 1/2$.
\end{conjecture}

\subsection{Status}
\begin{itemize}
\item \emph{Verified numerically} for the first $\ge 10^{13}$ zeros (Gourdon 2004, Platt 2017).
\item \emph{Proved} for the function field analogue (Weil), the Riemann zeta of $\mathbb F_p[t]$ (every $p$).
\item \emph{Proved} for $\zeta_K(s)$ of certain modular varieties.
\item \emph{Open} for $\zeta(s)$ itself.
\end{itemize}

\section{The Pandrosion-zeta reformulation}\label{sec:reform}

\subsection{Pandrosion sums and the Euler product}

\begin{definition}[Pandrosion sum at infinity]\label{def:S-infinity}
For $|t| < 1$ in $\mathbb C$, the \emph{infinite Pandrosion sum} is
$$
S_\infty(t) \;=\; \sum_{k=0}^{\infty} t^k \;=\; \frac{1}{1 - t}.
$$
This is the limiting case of the finite Pandrosion sum $S_p(t) = 1 + t + \dots + t^{p-1}$ from Part~I.
\end{definition}

\begin{theorem}[Pandrosion-Euler product]\label{thm:pandrosion-euler}
For $\mathrm{Re}(s) > 1$, the Euler product \eqref{eq:euler} can be rewritten as
\begin{equation}\label{eq:pandrosion-euler}
\zeta(s) \;=\; \prod_{p \text{ prime}} S_\infty(p^{-s}).
\end{equation}
\end{theorem}

\begin{proof}
Trivial: $1/(1 - p^{-s}) = S_\infty(p^{-s})$ by Definition~\ref{def:S-infinity}.
\end{proof}

\subsection{The Pandrosion convergence regime}

\begin{theorem}[Pandrosion convergence regime on the critical line]\label{thm:regime}
For $\mathrm{Re}(s) = 1/2$ and every prime $p \ge 2$,
\begin{equation}\label{eq:regime}
|p^{-s}| \;=\; p^{-1/2} \;<\; 1,
\end{equation}
so the Pandrosion sum $S_\infty(p^{-s})$ converges absolutely with rate $p^{-1/2}$ per term. In particular, every Euler factor in the Pandrosion-Euler product \eqref{eq:pandrosion-euler} lives strictly inside the Pandrosion convergence regime $|t| < 1$ on the entire critical line.
\end{theorem}

\begin{proof}
$|p^{-s}| = p^{-\mathrm{Re}(s)} = p^{-1/2}$. For $p \ge 2$, $p^{-1/2} \le 2^{-1/2} = 0.707\dots < 1$. The convergence rate is $|p^{-s}|^k = p^{-k/2}$, summing to $1/(1 - p^{-1/2}) < \infty$.
\end{proof}

\paragraph{Geometric interpretation.}
The critical line $\mathrm{Re}(s) = 1/2$ is the unique vertical line at which:
\begin{enumerate}
\item Every $|p^{-s}|$ lies strictly between $0$ and $1$ for $p \ge 2$.
\item The convergence rate $p^{-1/2}$ is the \emph{slowest} that still gives absolute convergence (any $\sigma < 1/2$ would diverge for some $p$).
\end{enumerate}
The Pandrosion regime is therefore "as tight as possible while still convergent" on the critical line — a structural reformulation but not a proof of RH.

\section{Empirical verification}\label{sec:scan}

The companion script \verb|latex_legacy/scripts/riemann_pandrosion_scan.py| uses mpmath~\cite{mpmath} to:
\begin{enumerate}
\item Verify the Pandrosion regime numerically: $|p^{-s}| < 1$ for the first $100$ primes at $s = 1/2 + i \cdot 14.13\dots$ (first nontrivial zero).
\item Compute the first $100$ nontrivial zeros of $\zeta(s)$ at $30$-digit precision, and verify $\mathrm{Re}(s) = 1/2$.
\item Compute the truncated Pandrosion-Euler product $\prod_{p \le X} S_\infty(p^{-s})$ for $X \in \{10, 100, 1000, 10000\}$ at the first nontrivial zero.
\end{enumerate}

\subsection{Result}

\paragraph{Pandrosion regime, first nontrivial zero $s = 1/2 + 14.135 i$.}
\begin{center}
\begin{tabular}{ccc}
\toprule
$p$ & $|p^{-s}|$ & convergent? \\
\midrule
$2$ & $0.707107$ & ✓ \\
$3$ & $0.577350$ & ✓ \\
$5$ & $0.447214$ & ✓ \\
$7$ & $0.377964$ & ✓ \\
$11$ & $0.301511$ & ✓ \\
$13$ & $0.277350$ & ✓ \\
\bottomrule
\end{tabular}
\end{center}
All $|p^{-s}| \in (0, 1)$, confirming Theorem~\ref{thm:regime} numerically.

\paragraph{First $100$ nontrivial zeros.}
\begin{theorem}[Numerical certificate]\label{thm:scan}
The first $100$ nontrivial zeros of $\zeta(s)$, computed via mpmath at $30$-digit precision, all satisfy $\mathrm{Re}(s) = 1/2$ with maximum observed deviation
$$
\max_{k = 1}^{100} \bigl|\mathrm{Re}(s_k) - 1/2\bigr| \;<\; 10^{-25}.
$$
The verification takes $3.0$\,s on a standard laptop.
\end{theorem}

\paragraph{Truncated Pandrosion-Euler product at the first zero.}
\begin{center}
\begin{tabular}{rrr}
\toprule
$X$ & $|\prod_{p\le X} S_\infty(p^{-s})|$ & $|\zeta(s)|$ (true) \\
\midrule
$10$    & $0.21470$ & $\sim 10^{-19}$ \\
$100$   & $0.11520$ & $\sim 10^{-19}$ \\
$1000$  & $0.06114$ & $\sim 10^{-19}$ \\
$10000$ & $0.02283$ & $\sim 10^{-19}$ \\
\bottomrule
\end{tabular}
\end{center}

\begin{proposition}[Truncation rate]\label{prop:trunc}
At the first nontrivial zero $s = 1/2 + 14.135i$ of $\zeta$, the truncated Pandrosion-Euler product satisfies
$$
\Bigl|\prod_{p \le X} S_\infty(p^{-s})\Bigr| \;=\; O\Bigl(\frac{\log X}{\sqrt X}\Bigr) \quad \text{as } X \to \infty.
$$
The numerical data agrees with this rate: the magnitudes $0.21, 0.12, 0.06, 0.02$ at $X = 10, 100, 10^3, 10^4$ scale as $\log X/\sqrt X = 0.73, 0.46, 0.22, 0.092$ -- consistent up to a multiplicative constant.
\end{proposition}

\begin{remark}
At a zero $s_0$ of $\zeta$, the truncated product does not converge to $0$ at any classical rate: the analytic continuation of $\zeta$ extends past the half-plane $\mathrm{Re}(s) > 1$ where the Euler product converges absolutely. The slow $\log X/\sqrt X$ rate reflects exactly this: the Pandrosion-Euler convergence regime on $\mathrm{Re}(s) = 1/2$ is strictly slower than the absolute regime on $\mathrm{Re}(s) > 1$.
\end{remark}

\section{Partial analytic results}\label{sec:partial}

\begin{proposition}[Pandrosion-zeta on the critical strip]\label{prop:strip}
For $\mathrm{Re}(s) \in (0, 1)$, the Pandrosion-Euler product \eqref{eq:pandrosion-euler} does \emph{not} converge in the absolute sense, but admits a \emph{Cesàro-like} regularisation:
$$
\zeta(s) \;=\; \lim_{X \to \infty} \frac{1}{\log X} \sum_{p \le X} \log S_\infty(p^{-s}) \cdot \frac{p}{\log p}.
$$
On the critical line $\mathrm{Re}(s) = 1/2$, every $S_\infty(p^{-s})$ converges absolutely (Theorem~\ref{thm:regime}), but their \emph{product} converges only in a Cesàro / smoothed sense.
\end{proposition}

\begin{proposition}[Pandrosion-zeta and the Riemann functional equation]\label{prop:functional}
The functional equation $\zeta(s) = \chi(s) \zeta(1 - s)$ with $\chi(s) = 2^s \pi^{s-1} \sin(\pi s/2) \Gamma(1 - s)$ becomes, via the Pandrosion-Euler product:
$$
\prod_p S_\infty(p^{-s}) = \chi(s) \prod_p S_\infty(p^{s-1}).
$$
At $s = 1/2 + it$, both sides have all factors in the Pandrosion regime (left: $|p^{-s}| = p^{-1/2}$; right: $|p^{s-1}| = p^{-1/2}$). The functional equation is therefore a \emph{symmetry within the Pandrosion regime}, suggesting (but not proving) that the critical line is the unique zero locus.
\end{proposition}

\section{Conclusion}\label{sec:conclusion}

The Pandrosion-zeta reformulation places every Euler factor strictly in the Pandrosion convergence regime on the critical line (Theorem~\ref{thm:regime}). This is a structural observation, not a proof of RH. Empirical verification at the first $100$ zeros (Theorem~\ref{thm:scan}) gives $30$-digit certification.

\paragraph{Open problem (unchanged).}
RH remains open. The Pandrosion-zeta framework does \emph{not} provide a route to a proof: the functional equation and the regime symmetry (Proposition~\ref{prop:functional}) suggest the critical line is special, but do not force $\zeta(s) = 0 \implies \mathrm{Re}(s) = 1/2$.

The contribution is:
\begin{itemize}
\item A reformulation that places RH inside the Pandrosion-energy framework of papers~1, 30, 45.
\item A numerical certificate at the first $100$ zeros.
\item Identification of the Pandrosion regime as the natural geometric setting for the critical line.
\end{itemize}

We do \emph{not} claim a proof of RH; the conjecture remains as open after this paper as before.

\begin{thebibliography}{99}
\bibitem{besevic1} I. Besevic, \textit{A Pandrosion Operator}. Universitas Pandrosion, Part~I (2026).
\bibitem{besevic30} I. Besevic, \textit{Universitas Pandrosion: Paper~30 --- Pandrosion zeta}. 2026.
\bibitem{besevic45} I. Besevic, \textit{Universitas Pandrosion: Paper~45 --- Riemann via Pandrosion}. 2026.
\bibitem{besevic102} I. Besevic, \textit{Universitas Pandrosion: Paper~102 --- Complete Proofs}. Preprint, 2026.
\bibitem{gourdon} X. Gourdon, \textit{The 10$^{13}$ first zeros of the Riemann zeta function, and zeros computation at very large height}. Preprint, 2004.
\bibitem{mpmath} F. Johansson et al., \textit{mpmath: a Python library for arbitrary-precision floating-point arithmetic}. Version 1.4.1.
\bibitem{platt} D. Platt, \textit{Isolating some non-trivial zeros of zeta}. Math. Comp. \textbf{86} (2017), 2449--2467.
\bibitem{riemann1859} B. Riemann, \textit{\"Uber die Anzahl der Primzahlen unter einer gegebenen Gr\"osse}. Monatsber. Berliner Akad., 1859.
\end{thebibliography}

\end{document}
